% --- CORRECCIÓN IMPORTANTE EN LA PRIMERA LÍNEA ---
% Agregamos "xcolor=table" para que funcione \rowcolor sin errores
\documentclass[aspectratio=169, 10pt, xcolor=table]{beamer}

% --- Configuración del Tema (estilo profesional como el ejemplo) ---
\usetheme{Madrid}
\usecolortheme{dolphin}

% --- Paquetes necesarios ---
\usepackage[utf8]{inputenc}
\usepackage[spanish]{babel}
\usepackage{amsmath, amssymb}
\usepackage{booktabs}
\usepackage{graphicx}
\usepackage{listings}
\usepackage{tikz}
\usetikzlibrary{arrows.meta, positioning, shapes.geometric}

% --- Ruta de Imágenes (crea la carpeta "imagenes" y coloca ahí tus archivos) ---
\graphicspath{{./imagenes/}}

% --- Estilo para código Python ---
\definecolor{codegreen}{rgb}{0,0.6,0}
\definecolor{codegray}{rgb}{0.5,0.5,0.5}
\definecolor{codepurple}{rgb}{0.58,0,0.82}
\definecolor{backcolour}{rgb}{0.95,0.95,0.92}
\lstset{
    language=Python,
    backgroundcolor=\color{backcolour},
    commentstyle=\color{codegreen},
    keywordstyle=\color{blue},
    numberstyle=\tiny\color{codegray},
    stringstyle=\color{codepurple},
    basicstyle=\ttfamily\scriptsize,
    breaklines=true,
    frame=single,
    showstringspaces=false
}

% --- Pie de página personalizado ---
\makeatletter

\setbeamertemplate{footline}{
  \leavevmode%
  \hbox{%
  \begin{beamercolorbox}[wd=.333333\paperwidth,ht=2.25ex,dp=1ex,center]{author in head/foot}%
    \usebeamerfont{author in head/foot}\insertshortauthor
  \end{beamercolorbox}%
  \begin{beamercolorbox}[wd=.333333\paperwidth,ht=2.25ex,dp=1ex,center]{title in head/foot}%
    \usebeamerfont{title in head/foot}Sesión 14
  \end{beamercolorbox}%
  \begin{beamercolorbox}[wd=.333333\paperwidth,ht=2.25ex,dp=1ex,right]{date in head/foot}%
    \usebeamerfont{date in head/foot}NLP \hfill \insertframenumber/\inserttotalframenumber \hspace*{0.3cm}
  \end{beamercolorbox}}%
  \vskip0pt%
}

\makeatother

% --- Datos de la portada ---
\title[SESIÓN 14 NLP]{\textbf{Sesión 14}}
\subtitle{Despliegue de Modelos NLP: MLOps, APIs, Contenerización y Monitoreo}
\author[Prof. C. Sánchez]{Carlos César Sánchez Coronel}
\institute{\textbf{NLP}}
\date{}

% Logo opcional (descomentar si tienes la imagen)
% \titlegraphic{\includegraphics[height=1.5cm]{logo.png}}

% --- Eliminar la barra de navegación (opcional) ---
\setbeamertemplate{navigation symbols}{}

% --- Índice al inicio de cada sección ---
\AtBeginSection[]{
  \begin{frame}
    \frametitle{Contenido}
    \tableofcontents[currentsection]
  \end{frame}
}

% ============================================================================
% INICIO DEL DOCUMENTO
% ============================================================================
\begin{document}

% --- Portada ---
\begin{frame}
    \titlepage
\end{frame}

% --- Índice general ---
\begin{frame}{Contenido}
    \tableofcontents
\end{frame}

% ============================================================================
% SECCIÓN 1: ARQUITECTURA DE UN CHATBOT EN PRODUCCIÓN
% ============================================================================
\section{Arquitectura de un Chatbot en Producción}

\begin{frame}{Arquitectura multicapa}
    \begin{columns}[t]
        \column{0.5\textwidth}
        \textbf{Frontend}
        \begin{itemize}
            \item Interfaz de usuario (web, app).
            \item Captura la consulta y muestra la respuesta.
        \end{itemize}
        \textbf{Backend}
        \begin{itemize}
            \item Lógica de negocio, gestión de sesiones.
            \item Se comunica con el modelo y BD.
        \end{itemize}
        \column{0.5\textwidth}
        \textbf{Base de datos}
        \begin{itemize}
            \item Historial de conversaciones, usuarios.
            \item Caché de respuestas (Redis).
        \end{itemize}
        \textbf{Modelo}
        \begin{itemize}
            \item Servicio de inferencia (API).
            \item Modelo serializado y optimizado.
        \end{itemize}
    \end{columns}
    % Basado en Pregunta 1
    % IMAGEN: Diagrama de arquitectura
\end{frame}

% ============================================================================
% SECCIÓN 2: SERIALIZACIÓN DE MODELOS
% ============================================================================
\section{Serialización de Modelos}

\begin{frame}{Formatos de serialización}
    \begin{table}
        \centering
        \begin{tabular}{l|c|c|c}
            \toprule
            \textbf{Formato} & \textbf{Ventajas} & \textbf{Desventajas} & \textbf{Uso típico} \\
            \midrule
            Pickle/Joblib & Simple, nativo Python & Inseguro, solo Python & Prototipos, notebooks \\
            ONNX & Interoperable, optimizado & Curva de aprendizaje & Producción, baja latencia \\
            TorchScript & Optimizado PyTorch & Solo PyTorch & Producción con C++ \\
            \bottomrule
        \end{tabular}
    \end{table}
    % Basado en Pregunta 2
\end{frame}

\begin{frame}[fragile]{Exportar a ONNX (PyTorch)}
    \begin{lstlisting}
import torch
import onnx

model = ...  # modelo entrenado
dummy_input = torch.randn(1, 3, 224, 224)

torch.onnx.export(
    model,
    dummy_input,
    "model.onnx",
    input_names=["input"],
    output_names=["output"],
    dynamic_axes={"input": {0: "batch_size"}, "output": {0: "batch_size"}}
)

# Verificar
onnx_model = onnx.load("model.onnx")
onnx.checker.check_model(onnx_model)
    \end{lstlisting}
    % Basado en Pregunta 2
\end{frame}

% ============================================================================
% SECCIÓN 3: CREACIÓN DE APIS CON FASTAPI
% ============================================================================
\section{Creación de APIs con FastAPI}

\begin{frame}{¿Por qué FastAPI?}
    \begin{itemize}
        \item \textbf{Asíncrono}: maneja alta concurrencia con \texttt{async/await}.
        \item \textbf{Rendimiento}: comparable a Node.js y Go.
        \item \textbf{Documentación automática}: Swagger UI, ReDoc.
        \item \textbf{Validación de datos}: Pydantic integrado.
        \item Ideal para servicios de ML.
    \end{itemize}
    % Basado en Pregunta 3
\end{frame}

\begin{frame}[fragile]{Endpoint básico en FastAPI}
    \begin{lstlisting}
from fastapi import FastAPI
from pydantic import BaseModel
import joblib

app = FastAPI()
model = joblib.load("model.pkl")

class Texto(BaseModel):
    texto: str

@app.post("/predict")
async def predict(data: Texto):
    pred = model.predict([data.texto])[0]
    return {"sentimiento": pred}
    \end{lstlisting}
\end{frame}

% ============================================================================
% SECCIÓN 4: CONTENERIZACIÓN CON DOCKER
% ============================================================================
\section{Contenerización con Docker}

\begin{frame}{¿Por qué Docker?}
    \begin{itemize}
        \item Empaqueta la aplicación y dependencias en una imagen.
        \item Garantiza el mismo entorno en desarrollo, pruebas y producción.
        \item Facilita el escalado y despliegue.
    \end{itemize}
    % Basado en Pregunta 4
\end{frame}

\begin{frame}[fragile]{Dockerfile para una API de NLP}
    \begin{lstlisting}
FROM python:3.9-slim

WORKDIR /app

COPY requirements.txt .
RUN pip install -r requirements.txt

COPY app.py .
COPY model.pkl .

CMD ["uvicorn", "app:app", "--host", "0.0.0.0", "--port", "80"]
    \end{lstlisting}
\end{frame}

\begin{frame}{Docker Compose}
    \begin{itemize}
        \item Define y levanta múltiples contenedores (API, base de datos, caché).
        \item Usa un archivo \texttt{docker-compose.yml}.
        \item Ejemplo: API + Redis para caché.
    \end{itemize}
\end{frame}

% ============================================================================
% SECCIÓN 5: MONITOREO Y MLOPS
% ============================================================================
\section{Monitoreo y MLOps}

\begin{frame}{Métricas críticas en producción}
    \begin{table}
        \centering
        \begin{tabular}{l|c}
            \toprule
            \textbf{Métrica} & \textbf{Descripción} \\
            \midrule
            Latencia & Tiempo de respuesta (p95 < 500ms) \\
            Throughput & Peticiones por segundo \\
            Data Drift & Cambios en distribución de entradas \\
            Concept Drift & Cambios en relación entrada-salida \\
            \bottomrule
        \end{tabular}
    \end{table}
    % Basado en PDF y Pregunta 5
\end{frame}

\begin{frame}{Herramientas de monitoreo}
    \begin{itemize}
        \item \textbf{Prometheus + Grafana}: métricas y dashboards.
        \item \textbf{ELK (Elasticsearch, Logstash, Kibana)}: centralización de logs.
        \item \textbf{Evidently AI}: detección de drift.
        \item \textbf{WhyLabs / Whylogs}: perfiles de datos.
    \end{itemize}
    % Basado en Pregunta 5
\end{frame}

\begin{frame}{Data Drift vs Concept Drift}
    \begin{columns}[t]
        \column{0.5\textwidth}
        \textbf{Data Drift}
        \begin{itemize}
            \item Cambios en las entradas.
            \item Ej: nuevas palabras, consultas más largas.
            \item Se detecta con PSI, test KS.
        \end{itemize}
        \column{0.5\textwidth}
        \textbf{Concept Drift}
        \begin{itemize}
            \item Cambios en la relación entrada-salida.
            \item Ej: nueva política cambia las respuestas correctas.
            \item Requiere reentrenamiento.
        \end{itemize}
    \end{columns}
    % Basado en Pregunta 6
\end{frame}

% ============================================================================
% SECCIÓN 6: VERSIONADO DE MODELOS Y DATOS
% ============================================================================
\section{Versionado de Modelos y Datos}

\begin{frame}{Herramientas de versionado}
    \begin{itemize}
        \item \textbf{DVC}: versiona datos y modelos, integrado con Git.
        \item \textbf{Hugging Face Hub}: repositorio de modelos.
        \item \textbf{MLflow}: registro de experimentos y modelos (Model Registry).
    \end{itemize}
    % Basado en Pregunta 7
\end{frame}

\begin{frame}{Estrategias de despliegue}
    \begin{itemize}
        \item \textbf{A/B testing}: comparar dos versiones en paralelo.
        \item \textbf{Canary deployment}: lanzar nueva versión a un % de usuarios.
        \item \textbf{Blue-Green}: dos entornos idénticos, cambiar el tráfico.
        \item \textbf{Shadow mode}: evaluar nueva versión sin impactar al usuario.
    \end{itemize}
    % Basado en Preguntas 12 y 13
\end{frame}

% ============================================================================
% SECCIÓN 7: ESCALABILIDAD Y RENDIMIENTO
% ============================================================================
\section{Escalabilidad y Rendimiento}

\begin{frame}{Escalado horizontal}
    \begin{itemize}
        \item Múltiples réplicas detrás de un balanceador de carga (Nginx, ALB).
        \item \textbf{Stateless}: cualquier réplica puede atender cualquier petición.
        \item Auto-escalado con Kubernetes (HPA) basado en CPU, memoria o colas.
    \end{itemize}
    % Basado en Preguntas 8 y 14
\end{frame}

\begin{frame}{Caché de respuestas}
    \begin{itemize}
        \item Usar Redis para respuestas de consultas frecuentes.
        \item Clave: hash del prompt (o embedding de la pregunta).
        \item TTL para expiración automática.
        \item Invalidación al actualizar el modelo.
    \end{itemize}
    % Basado en Pregunta 19
\end{frame}

% ============================================================================
% SECCIÓN 8: PRUEBAS DE CARGA
% ============================================================================
\section{Pruebas de Carga}

\begin{frame}{Herramientas y métricas}
    \begin{columns}[t]
        \column{0.5\textwidth}
        \textbf{Herramientas}
        \begin{itemize}
            \item Locust (Python)
            \item JMeter (Java)
            \item k6 (JavaScript)
        \end{itemize}
        \column{0.5\textwidth}
        \textbf{Métricas a observar}
        \begin{itemize}
            \item Latencia media y percentil 95
            \item Throughput (req/s)
            \item Tasa de errores (5xx)
        \end{itemize}
    \end{columns}
    \vspace{0.3cm}
    \begin{itemize}
        \item Determinar punto de quiebre aumentando carga gradualmente.
    \end{itemize}
    % Basado en Pregunta 11
\end{frame}

% ============================================================================
% SECCIÓN 9: SEGURIDAD Y PRIVACIDAD
% ============================================================================
\section{Seguridad y Privacidad}

\begin{frame}{Riesgos y mitigaciones}
    \begin{itemize}
        \item \textbf{Inyección de prompts}: validar y sanitizar entrada.
        \item \textbf{Extracción del modelo}: rate limiting, límites de consultas.
        \item \textbf{DoS}: autenticación (API keys, OAuth) y balanceo.
    \end{itemize}
    % Basado en Pregunta 9
\end{frame}

\begin{frame}{Cumplimiento GDPR}
    \begin{itemize}
        \item Almacenar logs de forma segura y con acceso restringido.
        \item Seudonimización: reemplazar identificadores por tokens.
        \item Derecho al olvido: implementar funciones de eliminación de datos.
        \item No retener datos personales más tiempo del necesario.
    \end{itemize}
    % Basado en Preguntas 10 y 18
\end{frame}

% ============================================================================
% SECCIÓN 10: ESTRATEGIAS DE DESPLIEGUE AVANZADAS
% ============================================================================
\section{Estrategias de Despliegue Avanzadas}

\begin{frame}{A/B Testing y Canary Deployment}
    \begin{block}{A/B Testing}
        \begin{itemize}
            \item Asignar usuarios aleatoriamente (hash de user\_id).
            \item Comparar métricas de negocio (CTR, conversión).
            \item Período suficiente para significancia estadística.
        \end{itemize}
    \end{block}
    \begin{block}{Canary}
        \begin{itemize}
            \item Enrutar un pequeño porcentaje de tráfico a la nueva versión.
            \item Monitorear errores y latencia.
            \item Aumentar gradualmente si es exitoso.
        \end{itemize}
    \end{block}
    % Basado en Preguntas 12 y 13
\end{frame}

\begin{frame}{CI/CD para ML (MLOps)}
    \begin{enumerate}
        \item Integración de código y datos.
        \item Entrenamiento automático con validación.
        \item Registro en MLflow.
        \item Despliegue canario / shadow.
        \item Monitoreo y alertas.
        \item Rollback automático si falla.
    \end{enumerate}
    Herramientas: GitHub Actions, Jenkins, Kubeflow.
    % Basado en Pregunta 12
\end{frame}

% ============================================================================
% SECCIÓN 11: PREPARACIÓN PARA ENTREVISTAS
% ============================================================================
\section{Preparación para Entrevistas}

\begin{frame}{Pregunta clave}
    \begin{block}{¿Cómo desplegarías un modelo NLP en producción?}
        \begin{enumerate}
            \item Serializar el modelo (ONNX para optimización).
            \item Crear una API con FastAPI (asíncrona, validación).
            \item Contenerizar con Docker.
            \item Orquestar con Kubernetes (auto-escalado, balanceo).
            \item Monitorear latencia, throughput y drift (Prometheus, Grafana).
            \item Versionar modelos (MLflow) y usar estrategias de despliegue (canary, A/B).
        \end{enumerate}
    \end{block}
    % Basado en PDF
\end{frame}

\begin{frame}{Preguntas típicas}
    \begin{itemize}
        \item ¿Qué arquitectura usarías para un chatbot?
        \item Compara formatos de serialización (Pickle, ONNX, TorchScript).
        \item ¿Cómo manejarías picos de tráfico?
        \item ¿Cómo detectarías data drift?
        \item ¿Qué herramientas de monitoreo conoces?
    \end{itemize}
\end{frame}

\begin{frame}{Preguntas avanzadas}
    \begin{itemize}
        \item ¿Cómo implementarías un A/B test para modelos?
        \item ¿Qué estrategias de caching propondrías para consultas frecuentes?
        \item ¿Cómo garantizarías el cumplimiento GDPR?
        \item ¿Qué es el shadow scoring y para qué sirve?
        \item ¿Cómo diseñarías un pipeline CI/CD para ML?
    \end{itemize}
\end{frame}

% ============================================================================
% SECCIÓN 12: CASO INTEGRADO - SISTEMA DE RECOMENDACIÓN
% ============================================================================
\section{Caso Integrado: Sistema de Recomendación}

\begin{frame}{Problema}
    \begin{itemize}
        \item Plataforma de e-commerce (Mercado Libre) necesita recomendar productos en tiempo real a millones de usuarios.
        \item Modelo basado en BERT (embeddings de productos y usuarios).
    \end{itemize}
    % Basado en Pregunta 20
\end{frame}

\begin{frame}{Arquitectura propuesta}
    \begin{itemize}
        \item \textbf{Serialización}: ONNX (baja latencia).
        \item \textbf{Backend}: FastAPI asíncrono.
        \item \textbf{Contenerización}: Docker.
        \item \textbf{Orquestación}: Kubernetes (auto-escalado horizontal).
        \item \textbf{Bases de datos}: Vectorial (FAISS/Pinecone) para productos, SQL para usuarios.
        \item \textbf{Caché}: Redis para consultas frecuentes.
    \end{itemize}
\end{frame}

\begin{frame}{Monitoreo y evolución}
    \begin{itemize}
        \item Métricas de sistema: latencia p95, throughput, errores (Prometheus/Grafana).
        \item Métricas de negocio: CTR, conversión, ingresos.
        \item Detección de concept drift: monitorear CTR; si baja, reentrenar.
        \item A/B testing para nuevas versiones.
        \item Seguridad: API keys, rate limiting, logs seudonimizados.
    \end{itemize}
\end{frame}

% ============================================================================
% SECCIÓN 13: RESUMEN
% ============================================================================
\section{Resumen}

\begin{frame}{Lo que hemos aprendido}
    \begin{exampleblock}{Conceptos clave}
        \begin{itemize}
            \item \textbf{Arquitectura}: frontend, backend, base de datos, modelo.
            \item \textbf{Serialización}: ONNX, TorchScript, Pickle.
            \item \textbf{APIs}: FastAPI (asíncrono, documentación automática).
            \item \textbf{Contenerización}: Docker, Docker Compose.
            \item \textbf{Monitoreo}: latencia, throughput, data drift, concept drift.
            \item \textbf{Herramientas}: Prometheus, Grafana, ELK, Evidently.
            \item \textbf{Versionado}: DVC, Hugging Face Hub, MLflow.
            \item \textbf{Escalado}: horizontal, caché (Redis).
            \item \textbf{Seguridad}: autenticación, rate limiting, GDPR.
            \item \textbf{Despliegue avanzado}: A/B testing, canary, CI/CD.
        \end{itemize}
    \end{exampleblock}
\end{frame}

% --- Diapositiva final ---
\begin{frame}
    \centering
    \Huge \textbf{¡Gracias!}

\end{frame}

\end{document}